\documentclass[11pt]{article}
\usepackage{fontspec}
\usepackage[french,english]{babel}
\usepackage{amsmath}
\usepackage{amssymb}
\usepackage{amsthm}
\usepackage{mathrsfs}
\usepackage{tikz-cd}
\usepackage{enumitem}
\usetikzlibrary{babel}
\tikzcdset{every diagram/.append style={baseline={([yshift=-.5ex]current bounding box.center)}}}

\newcommand{\printedpage}[1]{\marginpar{\footnotesize p.~#1}}
\newcommand{\msp}[1]{\marginpar{\footnotesize [#1]}}
\newcommand{\Hom}{\mathrm{Hom}}
\newcommand{\id}{\mathrm{id}}
\newcommand{\Winf}{\mathbf{W}_{\infty}}

\begin{document}

\printedpage{5}
% head: A. Grothendieck : Dérivateurs, Ch. XII, Édition des Universités de Montpellier II et Paris VII
% foot: Version 20 juin 2003

in $U'$, and hypotheses a), b) of page 1 ensure that the amalgamated sum $\varinjlim$ of this diagram in $U'^{\wedge}$ is $e_{U'}$, i.e. that
\[ p_!(e_U) = e_{U'}. \]
It follows that in the amalgamated sum, $X'$ considered in $(*)$, the unique element $t'$ of $T'$ is a final object of $X'$. We have seen that $U'$ is 1-connected, but not 2-connected if the enveloping abelian group of $G$ (i.e. $G/$(commutator subgroup) in the case where $G$ is a group) is not reduced to 1. The most economical example is obtained by taking $G = \{\pm 1\}$ the group with two elements. This is thus a case where in $(*)$ we have
\[ i \in W \ \ \text{(and even $i$ aspheric),} \qquad \text{but } i' \notin W, \]
and which has the advantage over the more sophisticated example of page 103 (prop. 10) of involving only finite categories -- and even very small ones!

Let us now proceed as in proposition 11, denoting by $U''$ the category obtained from $U'$ by adding to it a final
\msp{page 7}
object, and let
\[ p' : U' \longrightarrow U'' \]
be the inclusion, so that
\[
\begin{aligned}
p'_!(e_{U'}) \quad &= \quad (p'p)_!(e_U) = U' \quad && \text{as an object in } U''^{\wedge}\\
& && \text{(subpresheaf of the final presheaf } e_{U''}\text{).}
\end{aligned}
\]
Consider
\[
\begin{tikzcd}
U \arrow[r, hook, "i"] \arrow[d, "{p''=p'p}"'] \arrow[dr, phantom, "\begin{smallmatrix}\text{cocart.}\\ \text{and cart.}\end{smallmatrix}"] & X \arrow[d, "q''"'] & T \arrow[l, hook', dotted, "j"'] \arrow[d, dotted, "\text{iso}"] \\
U'' \arrow[r, hook, "i''"] & X'' & T''. \arrow[l, hook', dotted, "j''"']
\end{tikzcd}
\]
This time $U$, $X$ \textit{and} $U''$ are $\Winf$-aspheric (and even all three contractible), hence $i$, $p''$ are in $W$, yet $X''$ is not $\Winf$-aspheric (hence $i'', q'' \notin \Winf$), for in order for it to be so, i.e. for $i'' \in \Winf$, it would be necessary (prop. 7, cor. page 85 (\footnote[4]{not yet proved!})) that $i''$ be $\Winf$-aspheric, i.e. that $U''/t''$ be $\Winf$-aspheric. But by construction $U''/t'' \simeq U'$, which is \textit{not} $\Winf$-aspheric.

Here is an even simpler example of a functor $f : \Psi \longrightarrow X$ such that $f_!(e_{\Psi})$ is not $\Winf$-aspheric, i.e. of a diagram
\[
\begin{tikzcd}
& b \\
a \arrow[ur, "\beta_0"] \arrow[dr, "\gamma_0"'] & \\
& c
\end{tikzcd}
\]
\printedpage{6}
% head: A. Grothendieck : Dérivateurs, Ch. XII, Édition des Universités de Montpellier II et Paris VII
% foot: Version 20 juin 2003
\msp{page 8}
in a category $V$, such that the direct limit $L = b \amalg_a c$ \textit{in} $V^{\wedge}$ is a non-aspheric object of $V^{\wedge}$, i.e. $V/L$ is non-aspheric. I know that, for such a counterexample, $a \longrightarrow b$ and $a \longrightarrow c$ must not be monomorphisms. The idea is to take a universal example where neither of these arrows is \texttt{[a]} monomorphism, by taking
\msp{page 9}
\[
V = \left(
\begin{tikzcd}
& & b \\
0 \arrow[r, shift left, "\lambda"] \arrow[r, shift right, "\mu"'] & a \arrow[ur, "\beta"] \arrow[dr, "\gamma"'] & \\
& & c
\end{tikzcd}
\right)
\]
with
\[ \beta\lambda = \beta\mu \qquad \text{and} \qquad \gamma\lambda = \gamma\mu \]
the only two nonidentity arrows, besides the four arrows marked $\beta$, $\gamma$, $\lambda$, $\mu$.

We find that
\[ b \amalg_a c \xrightarrow{\ \sim\ } e_X \qquad \text{(amalgamated sum in } V^{\wedge}\text{).} \] % [sic] e_X -> e_V, the category considered here is V

Yet $V$ is not aspheric, nor even 1-connected: it is homotopic to the open subcategory $V_0 = (0 \overset{\lambda}{\underset{\mu}{\rightrightarrows}} a)$, for we have a retraction $V \xrightarrow{\ r\ } V_0$ (sending $b$ and $c$ to $a$), and an arrow $ir \longrightarrow \id_V$ (where $i : X_0 \hookrightarrow X$ is the inclusion). Thus $V$ has the same homotopy type \textit{as} $V_0$, which is a circle. So it is in our interest to take this $V$, where $\tilde{V}$ obtained from $V$ by adding a final element $e$, % [sic] i : X_0 \hookrightarrow X -> i : V_0 \hookrightarrow V, the inclusion of V_0 in V is meant
\[
\tilde{V} = \left(
\begin{tikzcd}
& & b \arrow[dr] & \\
0 \arrow[r, shift left] \arrow[r, shift right] & a \arrow[ur] \arrow[dr] & & e \\
& & c \arrow[ur] &
\end{tikzcd}
\right),
\]
as $U'$, for the counterexamples on pages 6 and 7.

\msp{page 10}
\subsection*{2.6. Counterexamples for total asphericity.}

I have finally convinced myself that
\[
\Psi = \left\{
\begin{tikzcd}
& b_0 \\
a_0 \arrow[ur] \arrow[dr] & \\
& c_0
\end{tikzcd}
\right\}
\]
\printedpage{7}
% head: A. Grothendieck : Dérivateurs, Ch. XII, Édition des Universités de Montpellier II et Paris VII
% foot: Version 20 juin 2003
is \textit{not} totally aspheric (for the localizer $\Winf$ -- in the present number, we shall assume $W = \Winf$), or what amounts to the same thing, that we can have a diagram of type $\Psi$
\[
\begin{tikzcd}
& b \\
a \arrow[ur, "\beta"] \arrow[dr, "\gamma"'] & \\
& c
\end{tikzcd}
\tag*{$(*)$}
\]
in a category $Z \in \mathrm{Ob}\,\mathrm{Cat}$, whose direct limit $b \amalg_a c$ in $Z^{\wedge}$ is the final presheaf $e_Z$, i.e. such that we have
\[ b(z) \amalg_{a(z)} c(z) \xrightarrow{\ \sim\ } \text{one-point set} \]
for every object $z$ of $Z$ (where $a(z) = \Hom_Z(a,z)$ \texttt{[rather} $\Hom_Z(z,a)$\texttt{]}) without $Z$ being aspheric for all that, nor even 1-connected (\footnote[5]{false -- it is necessarily 1-connected (\footnotemark[6]), but not 2-connected.}).\footnotetext[6]{cf. comment XII, p. 101.} The simplest example I have found, is the ``universal'' example of a diagram $(*)$ where neither one nor the other of the two arrows $a \longrightarrow b$, $a \longrightarrow c$ is \texttt{[a]} monomorphism (for if it is, then we can prove that $Z$ is indeed aspheric, as we shall see later on). It is the category
\[
Z = \left\{
\begin{tikzcd}
& & b_0 \\
0_0 \arrow[r, shift left, "\lambda_0"] \arrow[r, shift right, "\mu_0"'] & a_0 \arrow[ur, "\beta_0"] \arrow[dr, "\gamma_0"'] & \\
& & c_0
\end{tikzcd}
\right\}
\]
\msp{page 11}
obtained from $\Psi$ by adding to it an object $0_0$, connected to $a_0$ by two distinct arrows $\lambda_0$, $\mu_0$, such that
\[ \beta_0\lambda_0 = \beta_0\mu_0 \qquad \text{and} \qquad \gamma_0\lambda_0 = \gamma_0\mu_0, \]
these two composite arrows being, together with the four arrows $\beta_0$, $\gamma_0$, $\lambda_0$, $\mu_0$ marked on the diagram, the only (six) nonidentity arrows of $Z$. $\Psi$ immersed in $Z$ as a closed subcategory, with complement the one-point category reduced to $0_0$. (NB The closed immersion $\Psi \longrightarrow Z$ is not a Dwyer immersion, i.e. (as $Z$ is the smallest open subcategory of $Z$ containing $\Psi$) is not coaspheric, for the category $0\backslash\Psi$, cofibered over $\Psi$ with discrete fibers, of cardinality 2 at $a_0$ and 1 at $b_0$ and $c_0$, is the category
\[
\begin{tikzcd}
& \tilde{b}_0 \\
a'_0 \arrow[ur] \arrow[ddr] & \\
a''_0 \arrow[uur, crossing over] \arrow[dr] & \\
& \tilde{c}_0
\end{tikzcd}
\]
\printedpage{8}
% head: A. Grothendieck : Dérivateurs, Ch. XII, Édition des Universités de Montpellier II et Paris VII
% foot: Version 20 juin 2003
i.e.
\[
\begin{tikzcd}
& {} & \\
{} \arrow[ur] \arrow[dr] & & {} \arrow[ul] \arrow[dl] \\
& {} &
\end{tikzcd}
\quad ,
\]
which has the homotopy type of the circle.) Nor the open immersion $\{0_0\} \longrightarrow Z$, for $\{0_0\}/a_0$ is the discrete category of cardinality 2.

The inclusion $i : Z_0 \longrightarrow Z$ of
\[ Z_0 \ = \ (\ 0_0 \overset{\lambda_0}{\underset{\mu_0}{\rightrightarrows}} a_0\ ) \]
into $Z$ is a homotopism (\footnote[7]{false, there is no retraction at all! In fact, $Z$ is the \textit{suspension} of the $\rightrightarrows$, which has the homotopy type (topological) of the circle, and it therefore has the homotopy type of the 2-sphere.}) (for there exists an obvious retraction $r$ of $Z$ onto $Z_0$ (\footnotemark[7]), and an arrow $ir \longrightarrow \id_Z$), hence $Z$ is homotopic to $Z_0$, which has the homotopy type of the circle, hence is not 1-connected.

\msp{page 12}
On the other hand, the verification of
\[ b_0 \amalg_{a_0} c_0 \ = \ e_Z \qquad \text{in } Z^{\wedge} \]
is immediate.

We are going to use this arrow
\[ \Psi \longrightarrow Z \]
to construct interesting (and heartrending) examples of amalgamated sums
\[
\begin{tikzcd}
U \arrow[r, hook, "i"] \arrow[d, "p"'] & X \arrow[d, "q"'] \\
U' \arrow[r, hook, "i'"] & X',
\end{tikzcd}
\]
with $i$ an open immersion, such that we even have
\[ i :\ U \longrightarrow \overline{U} \qquad \text{is aspheric, hence in } W, \]
whereas nevertheless we shall have
\[ i' :\ U' \longrightarrow \overline{U'} \qquad \notin W, \]
and what is better (or worse),
\[ p \in W \qquad \text{but} \qquad q \notin W \]
\printedpage{9}
% head: A. Grothendieck : Dérivateurs, Ch. XII, Édition des Universités de Montpellier II et Paris VII
% foot: Version 20 juin 2003
(and possibly, at the same time, $i \in W$ but $i' \notin W$). Moreover (unlike the constructions of propositions 10, 11 (p. 101-107)), the categories considered may all be finite. I am particularly interested in the case where $X$ is a ``square''
\[ X = \Phi_n \times \Phi_m \qquad (n,m \geq 2), \]
and where $U$ is its ``boundary''
\[ U = \dot{X} \overset{\text{def}}{=} (\dot{\Phi}_n \times \Phi_m) \cup (\Phi_n \times \dot{\Phi}_m). \]

\msp{page 13}
\[
\begin{tikzcd}
0 \arrow[r] \arrow[d] & a \arrow[d] & a' \arrow[l] \arrow[r] \arrow[d] & \cdots \\
b \arrow[r] & t_1 & \cdots \arrow[l] & \\
b' \arrow[u] \arrow[r] \arrow[d] & {} \arrow[u] & \cdots \arrow[l] \arrow[u] & \\
\cdots & & &
\end{tikzcd}
\]

\textbf{Proposition.} \textit{Let $n,m \geq 2$, and consider the inclusion $U = \partial X \overset{i}{\hookrightarrow} X$ of the boundary of the ``square'' $X = \Phi_n \times \Phi_m$ into $X$. It is a $W$-aspheric open immersion, to be precise, the $U/t$ ($t \in \dot{U} = \bar{U}\backslash U$) are isomorphic to $e$ (one-point category) or to $\Phi =$ \begin{tikzcd}[ampersand replacement=\&, row sep=small, column sep=small] \& {} \arrow[dl] \\ {} \& \\ \& {} \arrow[ul] \end{tikzcd} (category with a final element), except for the four elements $t_i$ of the four ``corners'' of $\dot{U}$, where it is the category}
\[
\tilde{\Psi} = \left(
\begin{tikzcd}
& & a' \arrow[dl] \\
& a & \\
0 \arrow[ur] \arrow[dr] & & \\
& b & \\
& & b' \arrow[ul]
\end{tikzcd}
\right)
\]
\textit{having no final element. Thus $U \xrightarrow{\ i\ } X$ is not an open Dwyer immersion (but $U \longrightarrow X\backslash\{t_1,t_2,t_3,t_4\}$ is one \ldots). This being so, I say that $i : U \longrightarrow X$ is not a $W$-cofibration.}

\end{document}
